% Chapter 7: source-bound Spin(9)/Spin(8)-stabilizer geometry.
% Historical evidence: FCP S3L Oct 12; math Nov 7 strict runs 42426258689,
% 42429216538, 42452043377, 42498367504. No new Lean theorem is asserted.
\chapter[Source Spin(9) and the Ninth-Axis Stabilizer]{Source $\operatorname{Spin}(9)$, the Ninth-Axis Stabilizer, and the Geometry of the Cover}
\label{ch:spin-geometry}

\section{Why spatial symmetry must be built on the actual source}
The nine spatial labels in the original charge formula suggest the
$\operatorname{Spin}(9)$ double cover of the spatial rotation group. That
suggestion is not, by itself, a physical symmetry theorem. The matrix model
contains a real nine-dimensional bosonic vector index, a sixteen-dimensional
real spinor index, and a separate forty-eight-dimensional real Majorana
index. The spatial group must act compatibly on all of them; a rotation of
one index family cannot simply be declared a symmetry of the full closed
Hamiltonian.

The historical FCP S3L investigation exposed the difference. Its
strict-qualified kinetic free-spinor covariance concerned the \emph{actual}
source derivatives and spin matrices. The remaining cubic response involved
finite color sums, the order of the Majorana operators, and a complicated
composition of source identities. Smaller response lemmas qualified, but
the monolithic R4B/R4D final assembly encountered Lean recursion and
kernel-time limits. This was evidence about proof architecture, not a
counterexample to rotational covariance. The later \texttt{etblink/math}
program found more useful algebraic and geometric intermediate results
instead of simply repeating the same large application.

\begin{definition}[Original source spatial group and representations]
Write $V=\mathbb R^9$ with source quadratic form
$Q(v)=\sum_{i=0}^{8}v_i^2$, expressed in Lean as
\texttt{source9Quadratic}. Let
\[
 G_9:=\operatorname{spinGroup}(Q),\qquad
 \pi_9:G_9\longrightarrow \mathrm{SO}(V)
\]
denote the particular source Spin group and its \emph{source-defined}
vector representation. On the real sixteen-dimensional spinor index
there is a separately constructed homomorphism $D_{16}(g)$; on the
original real forty-eight Majorana labels its source realization is
\begin{equation}\label{eq:spin48-label-action}
 R_{48}(g)=E\bigl(I_3\otimes D_{16}(g)\bigr)E^{-1},
 \qquad g\in G_9,
\end{equation}
where $E$ is precisely the \texttt{pairedColorSpinEquiv} reindexing,
not a freely chosen relabelling of the Fock generators. In these equations
$R_{48}$ refers to an orthogonal \emph{label} representation, not
an operator on the complex Fock-vector space.
\end{definition}

The source group carries a Clifford/PBW topology fixed in the qualified
Lean construction. That detail matters: continuity in a conveniently
chosen abstract group model is not the same as continuity of the literal
homomorphisms used by the physical source. The later topological results
are valuable because they work in that specific source topology.

\section{Fixing the ninth axis without discarding the double cover}
Let $e_8$ denote the final standard coordinate vector, using the source
zero-based index $8$. There are two stabilizers that we must keep distinct:
\begin{equation}\label{eq:two-stabilizers}
 K:=\{g\in G_9:\pi_9(g)e_8=e_8\},\qquad
 H:=\{A\in\mathrm{SO}(9):Ae_8=e_8\}.
\end{equation}
Restriction of $\pi_9$ defines $p:K\to H$. The notation may suggest
a familiar pair $\operatorname{Spin}(8)\to\mathrm{SO}(8)$, but the
mathematical task is to establish exactly which source subgroup,
representation, topology, and covering map realize that familiar picture.
In particular, a group isomorphism unspecified on the vector eight-module
could introduce an unwanted triality identification.

\begin{proposition}[Exact source perpendicular geometry; strict-qualified]
The perpendicular source subspace
\[
 W=e_8^\perp\subset V
\]
has real dimension eight. The locus of source unit vectors perpendicular
to the ninth axis is the image of the radius-one sphere in $W$, hence
is path connected and compact. Moreover the source fixed-axis spin group
$K$ is the image of the continuous ordered \emph{eight-factor} Clifford
word map
\begin{equation}\label{eq:source-eight-word}
 \Phi:P:=\bigl\{u\in V:Q(u)=1,\ u_8=0\bigr\}^{8}
 \longrightarrow K.
\end{equation}
This is a statement about the \emph{entire} source $K$, including both
central signs, rather than merely a set of sample plane rotations.
\end{proposition}
\begin{proof}[Source-aligned mathematical explanation]
The exact transport \texttt{source9ToEuclidean} identifies the equation
$u_8=0$ with orthogonality to $e_8$. Its one-dimensional orthogonal
complement has dimension eight. Under that transport the conditions
$Q(u)=1$ and $u_8=0$ describe the unit sphere $S(W)$, which is
homeomorphic to $S^7$. It is path connected and compact.

An even Clifford word in perpendicular unit vectors fixes $e_8$:
reflections with perpendicular normals fix $e_8$, and the source
ordered reflection-to-Spin calculations preserve the multiplication
order. The more difficult converse uses the \emph{source} fixed-axis
Cartan factorization: every element of $K$ admits an ordered even
factorization with its unit normals in $W$, padded to eight factors.
Its exact central sign is handled within that source factorization.
Consequently the range of the continuous map $\Phi$ is exactly $K$.
This step is not justified solely by the familiar transitivity of
$\mathrm{SO}(9)$ on spheres; it uses the accepted original perpendicular
word theorem.
\end{proof}
\formalrecord{Source modules
\path{SpectralSourceNineAxisStabilizerWords.lean},
\path{SpectralSourceNineAxisStabilizerConnectedness.lean}; declarations
\texttt{source9Axis8Perpendicular\_finrank},
\texttt{source9Axis8EightUnitSpinWord\_range},
\path{source9Axis8EightUnitSpinWord_continuous_pbw}.
Strict-green path-connectedness gate
\href{https://github.com/etblink/math/actions/runs/42426258689}{\#42426258689}.}

\begin{theorem}[Compact and path-connected source stabilizer; strict-qualified]
\label{thm:source-K-compact-connected}
The actual subgroup $K$ is compact and path connected in its source PBW
topology. The map $\Phi:P\to K$ in \eqref{eq:source-eight-word} is a
continuous, surjective, closed quotient map. In particular, for any
topological space $Y$ and $f:K\to Y$,
\begin{equation}\label{eq:spin-quotient-continuity}
 f\text{ continuous}\quad\Longleftrightarrow\quad
 f\circ\Phi\text{ continuous}.
\end{equation}
\end{theorem}
\begin{proof}[Proof structure]
The parameter space $P\cong(S^7)^8$ is compact and path connected.
The preceding proposition gives continuity and surjectivity of $\Phi$.
Its image is thus compact and path connected, which proves the first
claim. As the original source PBW topology is Hausdorff, a continuous
map from compact $P$ to $K$ is closed. A continuous, surjective closed
map is a quotient map, giving \eqref{eq:spin-quotient-continuity}.
Notice that this argument gives a useful method for proving later
continuity statements on \emph{the real source group}: it is enough
to verify them after composition with an explicit eight-word map.
\end{proof}
\formalrecord{Strict-green quotient gate
\href{https://github.com/etblink/math/actions/runs/42429216538}{\#42429216538};
\path{SpectralSourceNineAxisStabilizerQuotient.lean},
named quotient-map and continuity equivalence declarations in atlas S-02.}

\section{The exact two-sheeted source covering}
The second part of the geometry controls the sign that a spatial
orthogonal matrix cannot see. In the source, the nontrivial negative
Spin center is denoted $z_9=-1$.

\begin{theorem}[Source ninth-axis covering; strict-qualified]
\label{thm:axis-covering}
The source projection $p:K\to H$ is surjective, continuous, and a
two-sheeted covering map. Its kernel and fibers are exactly
\begin{align}
 \ker p &=\{1,z_9\},\quad z_9\ne1,\label{eq:cover-kernel}\\
 p^{-1}(p(g))&=\{g,z_9g\},\qquad g\in K.\label{eq:cover-fibers}
\end{align}
The group $K$ is path connected and its target fixed-axis group has an
explicit source-coordinate equivalence
\begin{equation}\label{eq:axis-SO8}
 H\ \cong_{\mathrm{top.grp}}\ \mathrm{SO}(8).
\end{equation}
\end{theorem}
\begin{proof}[Proof structure and provenance]
The original Spin-vector action sends $z_9$ to the identity on $V$,
while the source two-sign fiber theorem proves that two spin elements
have the same full vector action only if they differ by $z_9$.
Applying this inside the fixed-axis subgroup gives
\eqref{eq:cover-kernel} and \eqref{eq:cover-fibers}.
The accepted source covering module combines the quotient property
of the projection with discreteness of this finite kernel to invoke
the covering-map criterion for a topological group homomorphism.
This proves the source covering statement rather than merely counting
preimages set-theoretically.

To identify $H$, restrict the first eight coordinates of a matrix in
$\mathrm{SO}(9)$ fixing $e_8$. Orthogonality implies that its last row
and last column are those of the identity, while the first $8\times8$
block belongs to $\mathrm{SO}(8)$. Conversely adjoining a final
$1\times1$ identity block yields a fixed-axis $\mathrm{SO}(9)$ matrix.
The source defines both maps, proves the group identities, and proves
both directions continuous, yielding \eqref{eq:axis-SO8} as an actual
\texttt{ContinuousMulEquiv}.
\end{proof}
\formalrecord{Original source files
\path{SpectralSourceNineAxisProjection.lean},
\path{SpectralSourceNineAxisCovering.lean},
\path{SpectralSourceNineAxisCoordinateEquivalence.lean}.
Strict-green covering gate
\href{https://github.com/etblink/math/actions/runs/42452043377}{\#42452043377}
and coordinate-equivalence gate
\href{https://github.com/etblink/math/actions/runs/42498367504}{\#42498367504}.
Principal declarations:
\texttt{source9Axis8ProjectionHom\_ker\_eq},
\path{source9Axis8ProjectionHom_isCoveringMap_pbw}, and
\path{source9Axis8CoordinateContinuousMulEquiv}.}

\section{Why this is not yet the branching theorem}
We now possess the correct original fixed-axis source cover and the
correct topological group on the orthogonal target. In classical compact
Lie-group theory, a connected two-sheeted spin covering of $\mathrm{SO}(8)$
is the setting in which $\operatorname{Spin}(8)$ representation theory
and the $B_4\downarrow D_4$ branching rule naturally appear. But the
\emph{specific} identification of the original $K$ with that standard
covering representation, its common maximal torus and weight lattice,
and the action on the original fermionic fiber have to be supplied at the
correct source types. An abstract group isomorphism is insufficient for
a theorem comparing the actual source vector representation with the
standard eight-dimensional vector module: $D_4$ triality can permute
its distinguished eight-dimensional representations.

This is also where signs and weights matter. The later analytical notes
corrected a proposed sign error: for the chosen active rotation,
$e_r+i e_s$ is a complexified \emph{basis vector} of weight $+1$,
while a coordinate function may transform with the opposite sign.
The four circles of a maximal torus do \emph{not} drop to three when
fixing one axis. The $B_4\to D_4$ restriction removes short roots,
not a torus circle. These corrections belong in the monograph because
an incorrect convention here would contaminate every later oscillator
and highest-weight bound.

\begin{openproblem}[Source-bound branching interface]
Construct the exact source vector-compatible covering equivalence
$K\to\operatorname{Spin}(8)$ (or an equivalent fully specified weight-
lattice interface), including both center signs, and prove the required
$\operatorname{Spin}(9)\downarrow\operatorname{Spin}(8)$ interlacing
property on the actual source representations. Then bind the relevant
representation to the original physical Fock-vector or operator-valued
fiber without changing its highest-weight convention.
\end{openproblem}

\section{Reader's checkpoint}
The original nine-dimensional Spin geometry is substantially realized:
we have actual source word parametrizations, connectedness, compactness,
a two-sheeted stabilizer covering, and the target $\mathrm{SO}(8)$
coordinate equivalence. The missing link is \emph{not} another generic
existence theorem for the Spin group. It is the precisely normalized
representation-theoretic identification needed by the original
physical-sector argument.
